\chapter{Physical and chemical background}\label{ch:background}
\section{The middle atmosphere and the model domain}
The mesosphere is a region in which solar radiation, chemical reactions and atmospheric motion jointly influence minor constituents. A concentration profile therefore need not remain fixed through a day, even when the major constituents vary slowly. This distinction is particularly relevant to ozone: its abundance is small compared with molecular oxygen, but its formation and photolysis connect the radiation field to reactive oxygen and hydrogen chemistry. The aeronomic context and the conventional distinction between major reservoirs and reactive minor species are described by \citet{brasseur2005}.

This thesis considers geometric heights from 50 to 100~km. The lower boundary includes the upper-stratospheric/lower-mesospheric transition; the upper boundary extends into the mesopause/lower-thermospheric region. The name ``mesospheric model'' is used as a concise description of this domain, not as a claim that every level belongs to a single sharply bounded atmospheric layer. Temperature, total neutral number density and reservoir concentrations are prescribed. Reactive concentrations evolve within this prescribed environment.

Two different uses of an atmospheric model must be distinguished. A background model supplies a temperature or density at a specified position and time. A chemical evolution model predicts how a minor-species state changes from a supplied initial condition. The latter cannot be obtained from geometry alone. Nor does prescribing time-dependent density imply that the calculation includes the material transport that would produce that density variation in the real atmosphere. This thesis intentionally separates these operations.

The calculation is a set of vertically resolved chemical boxes coupled by solar attenuation. There is no vertical advection, diffusion or horizontal transport of chemical constituents. Radiation can nevertheless couple different heights because a solar ray reaching one level passes through other parts of the column. Thus ``no transport'' means no constituent-transport terms in the differential equations; it does not mean that the radiation problem is local.

\section{Odd oxygen and ozone}
Atomic oxygen and ozone are linked through association and destruction. In the accepted reduced network, the elementary association event is
\begin{equation}\mathrm{O}+\oxygen+\mathrm{M}\longrightarrow\ozone+\mathrm{M},\label{eq:association}\end{equation}
where $\mathrm{M}$ denotes the collision partner represented by total neutral density. The corresponding event flux contains three concentrations. A coefficient with units of $\mathrm{cm^6\,molecule^{-2}\,s^{-1}}$ must therefore be multiplied by $[\mathrm{M}]$ exactly once. Ozone can also react with atomic oxygen, yielding molecular oxygen. These processes form part of the standard odd-oxygen framework \citep{brasseur2005,jpl18}.

Photolysis introduces an essential distinction between loss of a parent molecule and the branching of its products. Ozone absorbs solar radiation and may produce ground-state or excited products. A yield below unity for the excited channel does not reduce the total parent photolysis rate. If a fraction of the ozone photolysis events produces excited oxygen, the complementary events still destroy ozone. This distinction is explicit in the implemented reduced partition described in Chapter~\ref{ch:chemistry}.

Similarly, molecular-oxygen dissociation supplies two oxygen atoms per parent event. One retained channel supplies a ground-state atom and an excited O($^1$D) atom; the effective complementary channel supplies two ground-state atoms. Treating an excitation yield as if it were the entire dissociation frequency would violate this bookkeeping. The model uses total parent photolysis frequencies and separate excited-channel rates so that the reduced oxygen-atom account closes without pretending to retain every physical photoproduct state.

The term odd oxygen is useful because rapid interconversion can redistribute oxygen between atomic and ozone forms. The family is not simply conserved under all processes included here: reactions with prescribed reservoirs, dissociation of molecular oxygen and recombination alter its source and loss balance. In this work family diagnostics are checks on the event equations, not imposed conservation constraints that would override the chemistry. No total atmospheric oxygen inventory is evolved.

\section{Hydrogen radicals and peroxide}
The hydrogen-containing reactive species retained dynamically are H, OH, HO$_2$ and H$_2$O$_2$. They interact with ozone and atomic oxygen, while prescribed H$_2$O and H$_2$ provide sources or partners. OH and HO$_2$ can catalyse the conversion of reactive oxygen into molecular oxygen; their partition therefore matters even if their summed abundance appears less variable than either component. The reduced oxygen--hydrogen chemistry follows the historical model topology and evaluated coefficients identified in Chapter~\ref{ch:chemistry} \citep{li2020,jpl18}.

The family used in the accepted reduced formulation is
\begin{equation}R_H=[\mathrm{OH}]+[\mathrm{HO}_2].\label{eq:rhdef}\end{equation}
It does not include atomic H or H$_2$O$_2$. It is consequently a deliberately narrow diagnostic family rather than a complete hydrogen-atom inventory. For example, two HO$_2$ radicals can form one H$_2$O$_2$ molecule. This removes two members from $R_H$, although hydrogen remains within the wider set of hydrogen-bearing molecules.

The self-reaction illustrates the difference between event frequency and species consumption. With the accepted coefficient $k_{\mathrm{HO_2,HO_2}}$, its event flux is $k_{\mathrm{HO_2,HO_2}}[\mathrm{HO_2}]^2$. Each event consumes two HO$_2$ and produces one peroxide. Hence its contributions to $\dot R_H$ and $\dot{[\mathrm{H_2O_2}]}$ have coefficients $-2$ and $+1$, respectively. Adding a factor of one-half to this accepted event convention would change the model.

Peroxide photolysis releases two OH radicals in the retained network. OH reacting with peroxide produces HO$_2$, converting one radical form into another without changing $R_H$. These relationships make the family budget a useful diagnostic of implementation consistency. They also explain why a formally steady partition can fail when the radiation field changes: the reservoirs and radicals do not necessarily adapt on the same timescale.

\section{Excited oxygen and the infrared atmospheric band}
The notation $\Delta$ denotes the concentration of \singlet, the lowest singlet electronic state of molecular oxygen. Its emission near 1.27~$\mu$m provides a link between excited oxygen and mesospheric ozone. The connection has been used in ozone retrieval studies, including OSIRIS and SABER work \citep{li2020,mlynczak2007}. Those studies motivate the species of interest; their observational agreement does not validate the numerical results of this thesis.

Ozone photolysis in the Hartley region produces \singlet\ directly in the retained excited channel. The associated O($^1$D) can transfer energy through collisions with molecular oxygen. The model also includes molecular-oxygen excitation by the A, B and infrared atmospheric bands. Cascades and collisional transfer therefore provide sources of $\Delta$ in addition to the direct ozone channel. The need to account for sources other than ozone photolysis is established in the kinetic treatment of \citet{mlynczak1993}.

Two vibrational levels of the $b^1\Sigma_g^+$ state are represented: $B_0$ and $B_1$, corresponding to $v=0$ and $v=1$. Their concentrations, together with O($^1$D), are algebraic in the final model. In contrast, $\Delta$ is dynamic. Radiative decay and collisional quenching determine its loss frequency, and the finite lifetime allows its concentration to retain memory of earlier production.

An emission measurement and an excited-state concentration are not identical observables. In an optically simple local description, multiplying a concentration by a radiative coefficient produces a volume photon-production rate. A limb radiance additionally depends on viewing geometry, spectral response and transfer to the instrument \citep{mlynczak2007}. This thesis computes concentrations and chemical/radiative forcing fields. It does not implement a satellite retrieval or compare modelled limb radiances with observations.

\section{Solar photolysis and twilight}
A photolysis frequency is the number of parent dissociations per unit time and per parent molecule. It results from the incident photon field weighted by absorption cross sections and, where needed, product yields. Its units are $\mathrm{s^{-1}}$. Multiplication by the local parent concentration gives a volumetric event flux. Excitation frequencies for the oxygen bands have the same frequency units, although the parent is excited rather than dissociated.

The local solar field is attenuated by the column along the incoming ray. Near the horizon, a plane-parallel secant approximation becomes inappropriate: the ray can travel through curved shells or encounter the solid Earth. At a given solar zenith angle (SZA), higher levels may be illuminated while lower levels are in shadow. The onset of photochemistry therefore need not occur at one common time across the whole 50--100~km domain.

During dawn, production and loss coefficients can change on a timescale comparable to the adjustment of reactive species. An equilibrium inversion assumes that a measured excited-state abundance is in balance with instantaneous sources and sinks. A temporal initial-value calculation instead evolves the abundance from its prior state. The original project brief identifies this sunrise distinction as the motivation for replacing a stationary model. The present work examines that temporal forward problem; obtaining improved retrieved ozone is a subsequent application, not a demonstrated result here.

\section{Equilibrium, quasi-steady state and a temporal model}
Equilibrium of the whole retained network would require every evolved tendency to vanish. A quasi-steady-state approximation (QSSA) is narrower: selected fast intermediates are eliminated using instantaneous production and loss while other concentrations remain dynamic. These two operations should not be conflated. A model with three QSSA species can still be a temporal model with seven freely evolving concentrations.

The usefulness of a QSSA depends on both timescale separation and a physical algebraic solution. A small formal lifetime is not sufficient if the reduced equations require a negative concentration, or if their denominator vanishes while production remains positive. The temporal extension in this thesis provides concrete examples of these failures. The final reduction preserves only the three algebraic species whose positive denominators remain supported in the accepted domain.

The rest of the thesis separates background expectations from model definitions. Chapter~\ref{ch:chemistry} defines the accepted events, Chapter~\ref{ch:forcing} the prescribed environment and radiation, and Chapter~\ref{ch:temporal} the temporal reduction and initial conditions. Numerical validation and direct outputs are then presented independently of any atmospheric interpretation.
